\documentclass[12pt]{article}
\usepackage[T1]{fontenc}
\usepackage[utf8]{inputenc}
\usepackage{lmodern}
\usepackage{textcomp}
\usepackage{enumitem}
\usepackage{geometry}
\usepackage{graphicx}
\usepackage{amsmath, amssymb}
\geometry{margin=1in}
\begin{document}
\begin{center}
History - K-12 Level Assessment 7 \
Total Marks: 40 \
Answer all questions.
\end{center}

\section*{Multiple Choice (10 marks)}
\begin{enumerate}[label=\arabic*.]
\item Which dating technique is based on the patterns of tree-ring growth?
\begin{enumerate}[label=(\alph*)]
    \item radiocarbon dating
    \item dendrochronology
    \item paleomagnetic dating
    \item fission-track dating
\end{enumerate}
\item The record-keeping system used by the Inca in which a series of knotted strings were used as mnemonic devices is called:
\begin{enumerate}[label=(\alph*)]
    \item cuneiform.
    \item rachis.
    \item mit'a.
    \item khipu.
\end{enumerate}
\item What is one of the differences between simple foragers and complex foragers?
\begin{enumerate}[label=(\alph*)]
    \item Complex foragers employ irrigation technology.
    \item Complex foragers rely on a wider range of food sources.
    \item Complex foragers are more highly mobile.
    \item Complex foragers focus on a few highly productive resources.
\end{enumerate}
\item Between 1.4 million and 400,000 years ago, hand axe technology:
\begin{enumerate}[label=(\alph*)]
    \item steadily improved.
    \item steadily worsened.
    \item changed only slightly.
    \item evolved into a more sophisticated technology.
\end{enumerate}
\item The development of civilization in Mesopotamia is associated with the \_\_\_\_\_\_\_, during the \_\_\_\_\_\_.
\begin{enumerate}[label=(\alph*)]
    \item first writing; Amratian Period
    \item first cities; Uruk Period]
    \item first religious specialists; New Temple Period
    \item first permanent military; Nagada I Period
\end{enumerate}
\end{enumerate}

\section*{True / False (10 marks)}
\textit{Answer True or False}
\begin{enumerate}[label=\arabic*.]
\setcounter{enumi}{5}
\item True or False: Frere's discovery was located on the surface, above the remains of extinct animals.
\item True or False: The paleoanthropological perspective claims that the Earth is several billion years old and has undergone significant changes over time.
\item True or False: Ground-penetrating radar, proton magnetometer, and electrical resistivity are all remote-sensing devices used in archaeology.
\item True or False: Ancestral Puebloan culture in Chaco Canyon included around 75 substantial settlements and over 300 smaller communities.
\item True or False: Mohenjo-daro developed into a complex urban center during the Mature Harappan period, from 4,500 to 4,000 B.P.
\end{enumerate}

\section*{Long Form Response (20 marks)}
\textit{Answer each question with a clear and complete explanation.}
\begin{enumerate}[label=\arabic*.]
\setcounter{enumi}{10}
\item Discuss the significance of Homo habilis in the development of early human societies, particularly in relation to the creation and use of stone tools.
\item Discuss the significance of fire in human evolution, particularly focusing on how it enhanced the digestibility of food and its impact on health and society.
\end{enumerate}

\vfill
\noindent\textit{End of Paper}
\end{document}